\cleardoublepage
\newrefsection
\chapter{文献综述}

\section{背景介绍}
积雪深度（Snow Depth, SD）是刻画积雪储量与时空分布的重要指标，直接关系到水资源调配、融雪洪涝风险评估、生态与农牧业生产以及冰冻圈能量与物质交换过程。传统的雪深观测以地面站点、样线测量为主，虽然精度高，但存在观测效率低、空间覆盖稀疏、代表性不足等问题，难以满足大范围、多时相雪深动态监测需求。遥感观测能够提供大范围、周期性、多尺度雪深信息，已成为雪深监测的重要技术途径。

当前雪深遥感反演主要依赖被动微波、主动微波（尤其SAR）以及激光/雷达测高等数据源。被动微波具有全天候、长时间序列优势，是雪深与雪水当量反演的主力数据源，但空间分辨率较粗、易受混合像元与复杂地形影响，导致不确定性较大；SAR 具有更高空间分辨率，在小尺度雪深表征方面潜力突出，但反演机理复杂，积雪含水量、雪粒径与层结构等因素会显著影响后向散射；测高与激光雷达提供高精度高程观测，可通过“有雪/无雪地表高程差”提取雪深，为精细化雪深监测提供了新的可能。

\section{国内外研究现状}
由于北极的自然环境十分恶劣,获取大范围、高分辨率、准确的北极地区冰雪参数非常困难。为了满足积雪深度数据的使用需求,目前国内外针对积雪深度的获取途径开展了广泛的研究,包括冰面现场测量、浮标测量、船舶走航观测、机载航空测量和星载遥感反演。但上述五种方法的前四种只能获取局部的积雪深度信息,而星载遥感反演方法能够获取更广泛的空间分布数据。结合主动和被动微波遥感反演海冰表面积雪深度,尤其是利用合成孔径雷达数据进行相关研究,能够获取全北极尺度高分辨率的积雪深度信息,具有重要的研究意义和使用价值。本文将对国内外在现场积雪深度测量、被动微波反演积雪深度、主动微波反演积雪深度、积雪深度在海冰厚度反演和高分辨率船舶航线规划应用的研究进展进行综合叙述。
\subsection{研究方向及进展}
\subsubsection{基于现场观测数据的积雪深度研究}
作为环北极国家之一，前苏联从上个世纪三十年代开始在北极地区展开了一系列关于海冰表面积雪参数的科学研究。俄罗斯北极与南极研究所（Arctic and Antarctic Research Institute，AARI）的前身全联盟北极研究所（All-Union Arctic Institute）在1973年首次建立了位于西伯利亚海域与北极盆地区域的海冰漂流站，测量获得了长时间序列的降雪和海冰表面积雪数据，并将这些数据发布于美国冰雪数据中心（National Snow and Ice Data Center，NSIDC）的网站上（Colony et al.~\cite{colony1998}；Gerland and Haas~\cite{gerland2011}）。AARI共进行了一千余次北极科考，其中包括数十次高纬度科考行动，将共计约34个载人漂移冰站送到北极中央海域。基于1954年至1991年的AARI漂移冰站数据，Warren et al.~\cite{warren1999} 构建了整个北极海冰表面积雪深度模型（W99模型）。因为一年冰普遍较薄，不满足建立漂移冰站的条件，所以AARI漂移冰站获取的数据多为相对较厚的多年冰上覆盖的积雪深度数据，导致W99模型更适用于多年冰场景。距离W99模型的构建已过去31年，受全球变暖的影响，W99模型是否仍适用于当前北极海冰与积雪环境还有待确认。针对W99应用场景变化带来的准确性变化，相关学者在后续工作中对W99积雪深度模型进行了精度验证和校准。Kurtz and Farrell~\cite{kurtz_farrell2011} 和 Farrell et al.~\cite{farrell2012} 利用冰桥计划（Operation IceBridge，OIB）机载积雪深度观测值对W99模型反演的积雪深度进行了精度验证，证明了在一年冰环境下，准确的积雪深度值应为W99模型反演积雪深度数值的一半。基于上述研究，Laxon et al.~\cite{laxon2013} 和 Djepa et al.~\cite{djepa2014} 在反演海冰厚度时选择将多年冰和一年冰表面积雪深度进行区分，使用W99反演的积雪深度作为多年冰表面的积雪深度，而W99积雪深度的一半作为一年冰表面使用的积雪深度。

美国国家科学院于1974年提出利用浮标对整个北冰洋海域进行监测，标志着浮标在极地应用的开始。在随后的四年内，北极海洋浮标计划应运而生，并应用于海水表面温度、气压变化和海冰漂移轨迹的监测（郭井学等~\cite{guo2011}）。1991年，国际北极浮标计划（International Arctic Buoy Program，IABP）接续了北极海洋浮标计划，持续提供北极海冰的浮标观测数据。除了针对北极海冰研究的应用外，IABP实时监测数据为全球气候模式的研究也提供了重要的支持。

除了国际北极浮标计划（IABP）外，美国寒冷地区研究与工程实验室（the Cold Regions Research and Engineering Laboratory，CRREL）于1993年发布的物质平衡浮标（Ice Mass-Balance Buoy，IMB）数据也被广泛应用于北极海冰的研究（Perovich et al.~\cite{perovich2018}）。浮标数据采集利用卫星传输技术实现自动化，能够在无人值守情况下获取大量连续的高精度测点数据（Desch et al.~\cite{desch2017}）。在应用方面，Ricker et al.~\cite{ricker2014} 利用浮标获取的冰厚数据作为观测值，基于卫星测高反演了2013年秋季的北极海冰厚度，验证了积雪深度数据与海冰干舷存在正相关关系，结果表明积雪的积累会导致卫星测高反演海冰厚度相比实测数据存在高估。由于物质平衡浮标可以观测得到积雪、海冰、海水的空间位置，计算出积雪深度、海冰厚度与海冰温度，Gronfeldt~\cite{gronfeldt2016} 基于此数据建立了积雪深度反演海冰温度的模型，对模型反演海冰温度的精度进行了验证。Shalina and Sandven~\cite{shalina2018} 基于气象站定点测量的积雪深度数据与浮标数据，对北极东西伯利亚海和拉普捷夫海的海冰表面积雪深度进行分析，验证了东西伯利亚海的平均积雪深度可达15 cm，拉普捷夫海的平均积雪深度为9.8 cm，显著低于东西伯利亚海。Nicolaus et al.~\cite{nicolaus2018} 基于雪浮标（Snow Buoy，SB）数据证明北极的积雪深度与海冰厚度均存在逐年减少的情况。因北极自然环境恶劣，对积雪和海冰开展大范围、长时序的现场观测较为困难，浮标观测技术的应用能极大地缓解北极海冰积雪观测数据缺乏的状况，相较于站点的定点观测数据，浮标数据额外提供了海冰漂移的轨迹。

美国国家航空航天局（National Aeronautics and Space Administration，NASA）于2009年开始实施了为期六年的OIB实地观测计划，是有史以来对地球极地海冰进行的最大规模的空中调查。作为该计划的一部分，堪萨斯大学提供了雪雷达设备，将其搭载于机载平台上，获取了北极与南极海冰和冰川表面的积雪数据，并广泛应用于极地积雪深度的研究中。相关研究包括一年冰与多年冰表面积雪深度的差异、积雪深度的年际变化、积雪深度在海冰厚度反演中的作用等方面。Farrell et al.~\cite{farrell2012} 使用对应的站点积雪深度观测数据首次验证了OIB机载积雪深度的精度，论证了OIB积雪深度与对应的站点观测数据具有较好的一致性，平均偏差约为1 cm，在光滑的一年冰上二者相关系数为0.7。Brucker et al.~\cite{brucker2013} 使用OIB机载积雪深度作为观测数据，对高级微波扫描辐射计（Advanced Microwave Scanning Radiometer-Earth Observing System，AMSR-E）被动微波传感器反演的积雪深度进行验证，论证了AMSR-E反演的积雪深度在一年冰和多年冰场景下存在差异，应区分使用。Webster et al.~\cite{webster2014} 将1954—1991年AARI漂流冰站获取的积雪深度数据与2009年以后获取的OIB积雪深度数据进行了比较，发现随着年际变化，北极海冰表面积雪深度呈明显下降趋势。以波弗特海为例，积雪深度减少了约57\%。目前相关学者已经开发了多种基于OIB机载雪雷达数据反演积雪深度的算法，并公开发布了三种OIB积雪深度产品，分别为Kurtz et al.~\cite{kurtz2013} 发布的包含海冰干舷、积雪深度与海冰厚度参数的OIB四级产品和快速查看数据，以及Newman et al.~\cite{newman2014} 发布的基于Wavelet机载雪雷达的积雪深度产品。Kwok et al.~\cite{kwok_kurtz2017} 对上述三种积雪深度产品的精度进行了验证，结果表明Kurtz et al.~\cite{kurtz2013} 发布的产品往往低估了积雪深度，而Newman et al.~\cite{newman2014} 发布的积雪深度产品精度更高，与积雪深度实测数据相关系数可达0.72，但对10 cm以下的积雪深度存在高估。针对此情况，Kwok~\cite{kwok2018} 利用北极海冰干舷数据对OIB积雪深度数据进行了精度校正，结果表明仅在2014和2015年的实验校正效果较好，难以开展大范围的应用。针对OIB的机载雪雷达数据提出了一种新的获取积雪深度的方法，改善了OIB数据在低空与慢速条件下的观测精度。Rösel et al.~\cite{roesel2021} 基于P波段机载雪雷达数据，探讨了海冰表面积水对积雪深度估算的影响，提出了一种基于双极化雷达信号的方法来区分海冰表面积水与积雪，提高了反演积雪深度的准确性。

总的来说，基于机载雷达和地面观测的现场观测手段能够提供高精度的海冰表面积雪深度数据，但由于北极自然环境严苛，导致观测难度较大、经济成本较高、数据的时间和空间覆盖范围也存在局限性。因此，现场观测数据通常需要结合卫星遥感数据进行分析和验证，以期更全面地理解北极海冰表面积雪深度的时空变化规律。

\subsubsection{基于卫星被动微波遥感反演的积雪深度研究}
除了上文提到的AMSR-E之外，目前广泛应用于积雪深度反演的被动微波传感器还有三种，分别是特殊传感器微波成像仪（Special Sensor Microwave/Image，SSM/I）、专用传感器微波成像仪（Special Sensor Microwave Imager Sounder，SSMIS）和高级微波扫描辐射计2号（Advanced Microwave Scanning Radiometer 2，AMSR2）。基于星载传感器的被动微波遥感反演积雪深度研究始于1998年，Markus and Cavalieri~\cite{markus_cavalieri1998} 的研究证明被动微波遥感获取的海冰亮度温度与海冰表面积雪深度间存在较高的相关性，构建了南极威德尔海和别林斯高晋海的积雪深度观测值与SSM/I亮温数据的经验模型，开展了大范围的南极海冰表面积雪深度反演。Comiso et al.~\cite{comiso2003} 采用滑动平均的方法对被动微波遥感反演积雪深度经验模型的参数进行了校正与改进，使之适用于AMSR-E传感器，并反演了全北极范围的海冰表面积雪深度。Markus et al.~\cite{markus2006} 对上述模型进行了进一步校正，构建了Markus积雪深度反演模型，并对模型反演的积雪深度进行了精度评估，论证了Markus模型反演的积雪深度在多年冰上与观测数据偏差较大，更适用于一年冰场景。Rostosky et al.~\cite{rostosky2018} 对Markus模型在一年冰和多年冰表面积雪深度的估算进行了分析，证明Markus模型用于多年冰表面积雪深度反演存在明显的缺陷，需要引入更低频的被动微波数据来改善被动微波遥感反演多年冰积雪深度的效果。

随着Markus and Cavalieri~\cite{markus_cavalieri1998}、Comiso et al.~\cite{comiso2003}、Markus et al.~\cite{markus2006} 和 Rostosky et al.~\cite{rostosky2018} 的深入研究，不同的研究机构发布了不同传感器的被动微波积雪深度产品。NASA基于SSM/I被动微波亮温数据发布了北极海冰表面积雪深度产品，德国不莱梅大学（University of Bremen，UB）发布了SSMIS、AMSR-E和AMSR2被动微波积雪深度产品，NSIDC发布了AMSR-E被动微波积雪深度产品。然而，当前尚未进行汇总性的研究，以探讨和验证不同产品之间的差异和精度。在被动微波积雪深度产品的应用方面，Zwally et al.~\cite{zwally2008} 和 Kurtz et al.~\cite{kurtz2009} 利用AMSR-E积雪深度和卫星测高数据，基于海冰静力平衡模型反演了北极海冰厚度，论证了积雪深度的精度对海冰厚度的反演精度存在较大影响，在多年冰环境下存在较大误差。Markus and Cavalieri~\cite{markus_cavalieri2011} 和 Cavalieri et al.~\cite{cavalieri2012} 对基于Markus模型的多种被动微波遥感积雪深度进行了验证分析，证明随着年际变化，Markus模型反演一年冰表面积雪深度的精度存在下降趋势，精度降低的原因是随着全球变暖北极海冰减少，并存在由多年冰转化成一年冰的状况。针对AMSR-E和AMSR2同系列传感器，Lee et al.~\cite{lee2015} 定量分析了AMSR-E和AMSR2亮温数据反演的积雪深度差异，结果表明二者相关性可达0.9，平均偏差为1 cm。Castro-Morales et al.~\cite{castro_morales2017} 基于多种积雪深度观测数据对Markus模型和W99模型反演的积雪深度进行了精度验证，证明随着年际变化，这两种模型反演的积雪深度相对积雪深度观测数据都存在偏差增大的趋势，需对模型进行校正以获取新的海冰表面积雪深度数据。Lee et al.~\cite{lee2021} 综合使用AMSR-E、AMSR2和SMOS积雪深度产品，对2003—2020年间北极冬季积雪深度的空间分布差异和变化趋势进行了研究，证明了波弗特海与楚科奇海区域积雪深度的下降趋势显著高于其余区域，并探讨了造成这种差异的成因。Zhou et al.~\cite{zhou2021} 对比分析了两种被动微波遥感积雪深度数据与三种再分析数据的积雪深度，结果表明被动微波数据和再分析数据的积雪深度空间分布和年际变化趋势存在较大差异，并对差异产生的原因和积雪深度变化的驱动机制进行了研究。

综上所述，被动微波遥感能有效应用于获取大范围、长时序的北极海冰表面积雪深度，并可进一步结合海冰静力平衡方程反演北极海冰厚度。被动微波遥感反演的积雪深度能为北极地区大范围的海冰年际变化提供参考，但也存在被动微波穿透能力有限和空间分辨率较低的不足。

\subsubsection{基于卫星主动微波遥感反演的积雪深度研究}
主动微波遥感根据数据产品可分为非成像与成像遥感器两大类：非成像主动微波遥感器主要包括微波散射计和微波高度计；成像主动微波遥感器主要包括真实孔径雷达（Real Aperture Radar，RAR）和合成孔径雷达（Synthetic Aperture Radar，SAR）。

卫星测高反演海冰表面积雪深度的原理是利用高度计探测海冰—水面—雪面的位置信息差异，对积雪深度进行估算。其中，激光测高计探测位置为积雪表面，雷达测高计探测位置为海冰表面，同一位置的两者之差理论上即为海冰表面积雪深度。目前，Kwok et al.~\cite{kwok_kurtz2017} 利用冰、云和陆地高程2号卫星（ICESat-2）和2号冷卫星（CryoSat-2）所观测的干舷在同一分辨率栅格上的差值，对北极海域积雪深度进行了估算。Kwok and Markus~\cite{kwok_markus2018}、Kacimi and Kwok~\cite{kacimi_kwok2020} 和 Kacimi and Kwok~\cite{kacimi_kwok2022} 对CryoSat-2海冰干舷、ICESat-2总干舷与积雪深度之间的关系开展了一系列研究，验证了结合CryoSat-2和ICESat-2数据能够反演高精度的积雪深度与海冰厚度数据，为卫星测高海冰参数数据的广泛应用提供了有力支撑。此外，早期研究也发现，由船测数据获取的总干舷与积雪深度之间存在较好的相关性。Worby et al.~\cite{worby2011} 利用SIPEX 2007的船测数据探究了总干舷与海冰表面积雪深度之间的关系，发现通过海冰干舷阈值对冰型进行分类后，不同冰型的海冰干舷与积雪深度的线性拟合结果差异较大，但相关性均较高。Steer et al.~\cite{steer2016} 基于SIPEX、SIPEX II和ARISE获取的船测积雪深度与总干舷数据对上述结论进行了验证，发现不同船测数据在进行海冰分类时选取的海冰干舷阈值存在差异，且分类后积雪深度与总干舷数据的相关性优于分类前。

受限于雷达天线孔径，真实孔径雷达的方位向分辨率较低。在积雪参数反演方面，目前应用较为主流的是SAR。SAR积雪参数反演主要包括两类算法：一类是基于多频、多极化后向散射强度与积雪、土壤和植被覆盖等参数之间相关关系的经验算法；另一类是利用干涉SAR（InSAR）测量数据在观测积雪时产生的干涉相移，直接估算积雪深度。Shi and Dozier~\cite{shi_dozier2000a} 利用SIR-C/X-SAR获取的L、C和X波段VV和HH双极化SAR数据估算了雪水当量，该方法利用L波段估计雪密度及下垫面介电常数和粗糙度特性，在最小化地表后向散射信号影响的条件下，基于C和X波段估算雪深和粒径。Tedesco and Kim~\cite{tedesco_kim2006} 在大尺度遥感观测基础上发展了主被动微波遥感结合的雪深反演方法，相比单独使用主动或被动微波遥感，该方法反演结果的平均均方根误差（Root Mean Square Error，RMSE）降低了7.6\%，平均决定系数提高了6\%。Yu et al.~\cite{yu2022} 基于Sentinel-1数据，采用机器学习模型反演了美国德克萨斯州2021年的积雪深度，并结合气象数据与地面测量数据实现了积雪深度变化的预测。Hans et al.~\cite{hans2022} 基于哨兵1号（Sentinel-1）数据，利用后向散射系数的交叉极化组合，分别在100 m、500 m和1 km空间分辨率下反演了积雪深度，对于1.5–3 m范围内的积雪深度，其平均绝对误差约为观测值的20\%–30\%。

基于干涉相位的积雪深度反演目前主要应用于陆地。Guneriussen et al.~\cite{guneriussen2001} 最早利用InSAR反演雪水当量，论证了在Envisat SAR的C波段（波长5.62 cm）、入射角23°条件下，反演的积雪深度最大可达10.2 cm。针对新疆玛纳斯河流域的积雪特征，李晖等~\cite{lihui2014} 选用Envisat SAR数据，基于差分干涉技术反演了研究区积雪深度的空间分布。Mahmoodzada et al.~\cite{mahmoodzada2020} 基于Sentinel-1数据，采用差分干涉方法反演了阿富汗喜马拉雅山脉冬季和春季的积雪深度，并对反演结果进行了精度验证，进一步分析了反演过程中的不确定性。由于下垫面条件差异较大，且海冰存在生长和消融过程，与相对稳定的陆地差异显著，目前InSAR较少应用于海冰表面积雪深度的反演。

总体来看，在冰雪参数反演研究中，主动微波遥感的空间分辨率显著高于被动微波遥感，但针对极地地区的研究仍然较少，尤其是基于SAR数据的北极海冰表面积雪深度反演尚未开展深入研究。

\subsubsection{星载光子计数激光雷达}

星载光子计数激光雷达是一类能够基于单光子探测技术记录激光脉冲返回时间和空间位置的主动遥感系统，其代表性载荷是美国国家航空航天局（NASA）发射的ICESat-2卫星搭载的先进地形激光高度计（ATLAS）\citep{ICESat2_GSFC}。ICESat-2利用高重复频率的532\,nm激光脉冲和高灵敏度单光子探测器，以约10\,kHz 的发射频率沿轨道提供高精度地表高程测量，广泛用于冰盖、雪地和陆地地表变化监测\citep{Smith2019_ICESat2_HeightAlgorithm}。

在雪深反演领域，传统的激光雷达雪深估算通常依赖于积雪覆盖期与无雪期之间的地表高程差。具体方法是从ICESat-2的高度产品（如ATL06）中提取雪覆盖期的表面高度，再结合无雪期的参考高度或高分辨率数字高程模型（DEM）进行差分计算雪深\citep{DeschampsBerger2023_SnowDepth_ICESat2}。这种差分方法在中高纬度山区可以获得较高的雪深精度，但对辅助DEM的可用性和空间覆盖要求较高，在复杂地形和森林覆盖区的应用中仍存在误差问题\citep{DeschampsBerger2023_SnowDepth_ICESat2}。

随着光子计数激光雷达数据处理方法的发展，研究者提出了基于激光回波统计特征的“直接雪深反演”策略。Hu 等通过Monte Carlo激光雷达辐射传输模拟发现，雪层内部的多次散射导致激光回波的路径长度分布近似服从Gamma分布，并提出利用回波路径长度分布的统计矩作为雪深的反演量\citep{Lu2022_DerivingSnowDepth,sunmack2025_NeuralNetworkBased}。这些研究显示，光子返回路径的形态信息蕴含与雪深相关的物理信号，可在不依赖无雪期DEM的条件下实现雪深估算\citep{Lu2022_DerivingSnowDepth}。

基于多次散射路径统计的雪深反演方法在实际验证中也表现出良好的性能。例如，Lu 等的研究表明，在不同雪况条件下基于路径长度统计矩的方法能可靠估计雪深，并与参考数据存在较高的一致性\citep{Lu2022_DerivingSnowDepth}。此外，机器学习方法也开始与光子计数激光雷达数据结合，用于提高雪深估计的时空连续性，例如使用ICESat-2雪深生成训练数据，然后训练神经网络模型对其他传感器数据进行雪深预测\citep{sunmack2025_NeuralNetworkBased}。

总的来看，基于光子计数激光雷达的雪深反演研究主要分为两类思路：一类是基于雪期与无雪期高程差分的估算方法，另一类是基于激光回波统计特征的直接反演方法。后一种方法能够减少对外部数据的依赖，充分利用激光雷达回波信息，为高精度雪深监测提供了新的技术路径。

\subsection{存在问题与挑战}

近年来提出的基于星载光子计数激光雷达回波统计特征的直接雪深反演方法在一定程度上降低了对外部数据的依赖，为高精度雪深监测提供了新的技术路径。然而，该类方法目前仍面临多方面限制：一是缺乏针对雪深反演的系统物理理论模型，实际探测过程中激光雷达硬件参数（如激光能量分布、探测器在强信号条件下的非线性响应与饱和效应）以及环境因素（如云和气溶胶散射、地表坡度与粗糙度变化、太阳背景噪声等）均会对回波统计特征产生影响，从而增加反演不确定性；二是现有算法对复杂观测环境的适应性不足，在云雾影响显著或噪声水平较高的条件下反演性能下降明显；三是当前相关研究多基于有限区域或少量轨迹数据开展验证，其在大范围、多地形类型和复杂地表覆盖条件下的稳定性和适用性仍缺乏系统评估。

\subsection{研究展望}

针对上述基于星载光子计数激光雷达雪深反演中存在的关键问题，未来研究可围绕“仿真—反演—验证”三个层次系统推进相关方法的发展与完善。

在仿真层面，有必要构建面向雪深反演的光子计数激光雷达物理仿真模型，以系统刻画激光脉冲在雪层中的传播与散射过程。通过引入激光器参数（如脉冲能量分布、发射光斑形态）、探测器响应特性（如单光子探测效率、非线性响应与饱和效应）以及环境因素（如太阳背景噪声、云和气溶胶散射、地表坡度与粗糙度变化），可在仿真环境中分析不同因素对激光回波统计特征的影响机理，为雪深反演提供物理约束和理论依据。这类仿真研究有助于弥补当前雪深直接反演方法中物理模型不足的问题，并为算法设计提供可解释的参数空间。

在反演层面，应在物理机理约束的基础上发展更加稳健的雪深反演方法。一方面，可进一步挖掘激光回波多次散射路径长度分布及其统计特征与雪深之间的定量关系，构建具有明确物理含义的反演模型，从而降低对无雪期参考高程或外部DEM数据的依赖。另一方面，可结合机器学习与深度学习方法，在物理约束或先验条件的引导下建立数据驱动的反演模型，以增强算法在复杂观测环境和高噪声条件下的适应性和稳定性，实现对不同雪况和地形条件的有效刻画。

在验证层面，需要开展多尺度、多区域的系统性精度评估与适用性分析。目前相关研究多集中于局部区域或有限轨迹数据，未来应结合地面实测数据、机载激光雷达以及其他卫星遥感产品，在不同地形起伏、雪深等级和地表覆盖条件下对反演结果进行全面验证。同时，可通过跨传感器对比和时间序列分析，评估反演方法在大范围应用中的稳定性和一致性，从而为构建高精度、时空连续的雪深产品奠定基础。

通过仿真研究揭示物理机理、反演方法提升算法鲁棒性以及大范围验证评估反演精度与适用性，将有助于推动星载光子计数激光雷达雪深反演技术从方法探索向业务化应用迈进。

\newpage
\begingroup
    \linespreadsingle{}
    \printbibliography[title={参考文献}]
\endgroup
